\subsection{From a given power to the pulse area}

The pulse area is the only quantity the excitation depends on, and it follows
from the optical power in three steps: fix the scale of the intensity, convert
intensity into a Rabi frequency, integrate over the pulse. This section keeps
the three normalisations apart, because only one of them is physics.

\paragraph{What is normalised, and what is not.}
\begin{enumerate}
\item \emph{The mode.} $u(\mathbf 0)=1$, so the simulated intensity
      $\tilde I(\mathbf r,t)=|E|^2$ is in arbitrary units. This choice is free
      and drops out of every physical result.
\item \emph{The power.} The optical power $P$ in the profile is the single
      physical input. It fixes the scale factor $s$ in
      $I(\mathbf r,t)=s\,\tilde I(\mathbf r,t)$.
\item \emph{The reference.} A reference intensity $I_\mathrm{ref}$, an
      effective area $A_\mathrm{eff}$ and a nominal Rabi frequency
      $f_\mathrm{Rabi}$ are bookkeeping. They make the numbers readable and
      cancel from the pulse area.
\end{enumerate}

\paragraph{Step 1: the scale from the power.}
The power carried by the pattern is the area integral of the time average,
\begin{equation}
P=\int\!\langle I(\mathbf r)\rangle_t\,\mathrm d^2r
 =s\int\!\langle\tilde I(\mathbf r)\rangle_t\,\mathrm d^2r
 \equiv s\,\tilde P_\mathrm{prof},
\qquad\Longrightarrow\qquad
s=\frac{P}{\tilde P_\mathrm{prof}} .
\end{equation}
$\tilde P_\mathrm{prof}$ is known in closed form. A single spot of unit
amplitude carries
\begin{equation}
P_1=\int\!|u|^2\,\mathrm d^2r=
\begin{cases}
\dfrac{\pi w_0^2}{2}, & \text{Gauss } u=e^{-r^2/w_0^2},\\[2ex]
\dfrac{4\pi}{k^2}, & \text{Airy } u=\dfrac{2J_1(kr)}{kr},\quad
k=\dfrac{3.8317}{1.4830\,w_0},
\end{cases}
\end{equation}
and, since the cross terms of non-degenerate pairs average to zero in time,
\begin{equation}
\tilde P_\mathrm{prof}
=\underbrace{\sum_s a_s\,P_1}_{\text{incoherent}}
+\underbrace{2\!\!\sum_{\substack{s<s'\\ k_s=k_{s'}}}\!\!\sqrt{a_sa_{s'}}
\,\cos(\phi_s-\phi_{s'})\,
\int\! u(\mathbf r-\mathbf c_s)\,u(\mathbf r-\mathbf c_{s'})\,\mathrm d^2r}
_{\text{frequency degenerate pairs}} ,
\end{equation}
with the overlap integral, again analytic,
\begin{equation}
\int\! u(\mathbf r-\mathbf c_s)u(\mathbf r-\mathbf c_{s'})\,\mathrm d^2r
=P_1\cdot
\begin{cases}
e^{-d^2/2w_0^2}, & \text{Gauss},\\[1ex]
\dfrac{2J_1(kd)}{kd}=u(d), & \text{Airy},
\end{cases}
\qquad d=|\mathbf c_s-\mathbf c_{s'}| .
\end{equation}
The Airy result follows because the mode is the Fourier transform of a disc and
the autocorrelation of a disc is the disc again. The sum runs over the
intensity weights $a_s$, not $a_s^2$: the field carries $\sqrt{a_s}$. For
$3\times4$ at all phases zero the degenerate corner pair contributes $+1.1\,\%$.
The integral is never taken from the computation grid, which truncates about
$5\,\%$ of the Airy power at the usual margin.

\paragraph{Step 2: intensity to Rabi frequency.}
The two Raman legs, of intensities $\beta I$ and $(1-\beta)I$, drive
$|F{=}2,m\rangle\leftrightarrow|F{=}3,m\rangle$ via the $5P$ manifold. The
single-photon Rabi frequency of leg $i$ on the transition to an excited state
$e$ is $\Omega_{i,e}=d_e E_i/\hbar$, and with $I_i=\tfrac12\varepsilon_0cE_i^2$
\begin{equation}
\Omega_{i,e}^2=\frac{2\,d_e^2}{\varepsilon_0 c\,\hbar^2}\,I_i
\equiv d_e^2\,\kappa\,I_i,\qquad
\kappa=\frac{2}{\varepsilon_0 c\,\hbar^2} .
\end{equation}
For $|\Delta_e|\gg\Omega,\Gamma$ the excited states are eliminated
adiabatically. Each of them contributes one Raman path, and the paths add
coherently:
\begin{equation}
\Omega=\sum_e\frac{\Omega_{1,e}\,\Omega_{2,e}}{2\Delta_e}
=\underbrace{\kappa\sqrt{\beta(1-\beta)}
\sum_{e}\frac{d_{2e}\,d_{3e}}{2\Delta_e}}_{\textstyle C_\mathrm{rabi}}\;I ,
\end{equation}
The index $e$ labels an \emph{intermediate state}, i.e.\ a hyperfine level of
the $5P$ manifold,
\begin{equation}
|e\rangle=|5P_{J'},F',m'\rangle,\qquad
J'=\tfrac12\ (D_1,\ F'=2,3),\quad J'=\tfrac32\ (D_2,\ F'=1,2,3,4),
\end{equation}
so six terms for $\sigma^+$ light, where $m'=m+1$ is fixed by the polarisation.
$d_{2e}=\langle e|\hat d_{q}|F{=}2,m\rangle$ and
$d_{3e}=\langle e|\hat d_{q}|F{=}3,m\rangle$ are the \emph{signed} dipole matrix
elements, and
\begin{equation}
\Delta_e=\Delta+\left(\nu_{D_1}-\nu_\mathrm{line}\right)-\Delta_\mathrm{hfs}(F')
\end{equation}
is the detuning of the two-photon resonant pair from that level, measured from
the $5P_{1/2}$ centroid. The signs matter: the $F'$ contributions partly cancel,
and dropping them is wrong by tens of percent as soon as $|\Delta|$ approaches
the excited-state hyperfine splitting.

The formula above is evaluated with atomic data from the ARC package: the
signed matrix elements come from \texttt{getDipoleMatrixElementHFS}, the
hyperfine shifts from \texttt{getHFSCoefficients}/\texttt{getHFSEnergyShift},
the line centroids from \texttt{getTransitionFrequency} and the linewidths from
\texttt{getStateLifetime}. Only the elimination itself is coded by hand, in the method \texttt{coeff()}
of the module \texttt{rb85\_raman}. The same elimination gives the light shift of
each ground state and the photon scattering rate,
\begin{equation}
\delta_F=\kappa\sum_e d_{Fe}^2
\left[\frac{\beta}{4\Delta_e^{(1)}}+\frac{1-\beta}{4\Delta_e^{(2)}}\right]I,
\qquad
\Gamma_\mathrm{sc}=\kappa\sum_e\Gamma_e\,d_{Fe}^2
\left[\frac{\beta}{4(\Delta_e^{(1)})^2}+\dots\right]I ,
\end{equation}
so that all three are linear in $I$:
\begin{equation}
\Omega=C_\mathrm{rabi}I,\qquad
\delta_\mathrm{LS}=\delta_3-\delta_2\equiv C_\mathrm{shift}I,\qquad
\Gamma_\mathrm{sc}=C_\mathrm{scatter}I .
\end{equation}
Their ratio $\eta=\delta_\mathrm{LS}/\Omega$ is therefore a pure number,
independent of intensity, position and time. For $\sigma^+/\sigma^+$
co-propagating light on the clock states of Rb-85 and $\beta=\tfrac12$ one finds
$C_\mathrm{rabi}=11.66\,\mathrm{rad\,s^{-1}}/(\mathrm{W/m^2})$ and
$\eta=-0.061$ at $\Delta=+50$\,GHz, and $71.98$ and $-0.442$ at
$\Delta=-8$\,GHz. To leading order
$\eta\simeq-\omega_\mathrm{HF}/\Delta$, since the differential shift is the
difference of two denominators separated by the ground hyperfine splitting.

\paragraph{Step 3: the pulse area, without any reference.}
A rectangular pulse of length $T_p$ starting at $t_0$ accumulates the area
$\theta=\int\Omega\,\mathrm dt$. Three substitutions turn this into the working
formula. First the atomic law of Step 2, then the scale of Step 1, and finally
$s$ itself:
\begin{align}
\theta(\mathbf r,t_0)
&=\int_{t_0}^{t_0+T_p}\!\!\Omega(\mathbf r,t)\,\mathrm dt
&&\text{definition}\\
&=C_\mathrm{rabi}\int_{t_0}^{t_0+T_p}\!\! I(\mathbf r,t)\,\mathrm dt
&&\Omega=C_\mathrm{rabi}I\\
&=C_\mathrm{rabi}\,s\int_{t_0}^{t_0+T_p}\!\!\tilde I(\mathbf r,t)\,\mathrm dt
&&I=s\,\tilde I\\
&=C_\mathrm{rabi}\,\frac{P}{\tilde P_\mathrm{prof}}
  \int_{t_0}^{t_0+T_p}\!\!\tilde I(\mathbf r,t)\,\mathrm dt
&&s=P/\tilde P_\mathrm{prof} .
\end{align}
With $\tilde P_\mathrm{prof}=\int\langle\tilde I\rangle_t\,\mathrm d^2r$ this is
\begin{equation}
\boxed{\;
\theta(\mathbf r,t_0)
=C_\mathrm{rabi}\,P\;
\frac{\displaystyle\int_{t_0}^{t_0+T_p}\!\tilde I(\mathbf r,t)\,\mathrm dt}
     {\displaystyle\int\!\langle\tilde I\rangle_t\,\mathrm d^2r}\; }
\end{equation}
This is the whole physics. Only three things enter: the atomic coefficient, the
power, and the shape of the pattern, the latter as the ratio of a time integral
to an area integral. The arbitrary mode normalisation cancels in that ratio,
and no region, no $A_\mathrm{eff}$ and no $f_\mathrm{Rabi}$ appear. Units:
$[C_\mathrm{rabi}]\,\mathrm{W}\cdot\mathrm{s}/\mathrm{m^2}=\mathrm{rad}$.

\paragraph{Why $f_\mathrm{Rabi}$ may still be used.}
It is not an extra assumption but a name for the scale. Pick any reference
intensity $I_\mathrm{ref}$ -- in this work the time average over the evaluation
region -- and define
\begin{equation}
A_\mathrm{eff}=\frac{\int\langle I\rangle_t\,\mathrm d^2r}{I_\mathrm{ref}}
=\frac{\tilde P_\mathrm{prof}}{\langle\tilde I\rangle_\mathrm{ref}},
\qquad
I_\mathrm{ref}=\frac{P}{A_\mathrm{eff}},
\qquad
2\pi f_\mathrm{Rabi}\equiv C_\mathrm{rabi}I_\mathrm{ref} .
\end{equation}
Substituting $P=I_\mathrm{ref}A_\mathrm{eff}$ into the boxed equation gives
\begin{equation}
\theta(\mathbf r,t_0)
=2\pi f_\mathrm{Rabi}\,
\frac{1}{\langle\tilde I\rangle_\mathrm{ref}}
\int_{t_0}^{t_0+T_p}\!\tilde I(\mathbf r,t)\,\mathrm dt ,
\end{equation}
which is the same number: $A_\mathrm{eff}$ cancels against
$\tilde P_\mathrm{prof}$. Choosing a different region changes $I_\mathrm{ref}$,
$A_\mathrm{eff}$ and $f_\mathrm{Rabi}$ together and leaves $\theta$ untouched
(checked numerically to $10^{-14}$). $f_\mathrm{Rabi}$ is the Rabi frequency one
would measure by calibrating on the time averaged profile, without a trigger,
and $\Omega(\mathbf r,t)=C_\mathrm{rabi}I(\mathbf r,t)$ is what the atom
actually sees -- a quantity that varies in space and, through the beating,
strongly in time.

\paragraph{Where the sinc comes from.}
The time integral in the box is the only part that still depends on the
beating, and it can be done term by term. Inserting the Fourier series
$\tilde I=\sum_d D_d(\mathbf r)\,e^{i2\pi df_0t}$, each order contributes
\begin{equation}
\int_{t_0}^{t_0+T_p}\!\! e^{i2\pi df_0t}\,\mathrm dt
=\frac{e^{i2\pi df_0(t_0+T_p)}-e^{i2\pi df_0t_0}}{i\,2\pi df_0} .
\end{equation}
Pulling out the phase of the window \emph{centre}, $e^{i2\pi
df_0(t_0+T_p/2)}$, makes the bracket symmetric,
\begin{equation}
=e^{i2\pi df_0\left(t_0+T_p/2\right)}\,
\frac{e^{i\pi df_0T_p}-e^{-i\pi df_0T_p}}{i\,2\pi df_0}
=e^{i2\pi df_0\left(t_0+T_p/2\right)}\,
\frac{\sin\!\left(\pi df_0T_p\right)}{\pi df_0} ,
\end{equation}
where $e^{ix}-e^{-ix}=2i\sin x$ was used. Writing the last factor as $T_p$
times $\operatorname{sinc}(z)=\sin(\pi z)/(\pi z)$ gives
\begin{equation}
\int_{t_0}^{t_0+T_p}\!\tilde I\,\mathrm dt
=T_p\sum_d D_d(\mathbf r)\,\operatorname{sinc}(d f_0T_p)\,
e^{i2\pi df_0\left(t_0+T_p/2\right)} .
\end{equation}
The three factors read directly: $T_p$ is the length of the window, the
exponential says that only the \emph{centre} of the window matters for the
phase, and the sinc is the Fourier transform of the window itself, evaluated at
the beat frequency $df_0$. Integrating over a window of length $T_p$ is a
convolution with a boxcar, and in frequency space a convolution is a
multiplication -- with the transform of the box. The $d=0$ term is the limit
$\operatorname{sinc}(0)=1$: the time average always survives, whatever the
pulse length. Because $D_{-d}=D_d^*$ the sum is real,
\begin{equation}
\begin{split}
\int_{t_0}^{t_0+T_p}\!\tilde I\,\mathrm dt
=T_p\Bigl[\,&D_0(\mathbf r)\\
&+2\sum_{d>0}\left|D_d(\mathbf r)\right|
\operatorname{sinc}(df_0T_p)\,
\cos\!\left(2\pi df_0\left(t_0+\tfrac{T_p}{2}\right)+\arg D_d\right)\Bigr].
\end{split}
\end{equation}
This is the same expression as a camera exposure of duration $T_p$; a pulse and
an exposure differ only in their length. Two limits:
\begin{itemize}
\item $T_p\gg T_0$: every $d\neq0$ is damped away and
      $\theta\to C_\mathrm{rabi}\,s\,\langle\tilde I\rangle_t\,T_p$. The pulse
      sees the time average.
\item $T_p\ll T_0$: $\operatorname{sinc}\to1$. The pulse does not average the
      beating, it \emph{samples} it, and $\theta$ depends on the trigger
      instant $t_0$ -- on the time scale of the fast beats, not of $T_0$.
\end{itemize}

\paragraph{Averages over the atom are integrals.}
$\theta(\mathbf r)$ is a field. What the atom experiences is the average over
its position probability density, a proper integral,
\begin{gather}
\langle\theta\rangle_W=\int\! W(\mathbf r)\,\theta(\mathbf r)\,\mathrm d^2r,
\qquad \int\! W(\mathbf r)\,\mathrm d^2r=1,\\
W(\mathbf r)=\frac{1}{2\pi\sigma^2}
\exp\!\left[-\frac{(\mathbf r-\mathbf r_0)^2}{2\sigma^2}\right],
\end{gather}
with the thermal width
$\sigma^2=\frac{\hbar}{2m\omega_r}\coth\frac{\hbar\omega_r}{2k_BT}$
($\sigma=108$\,nm at $\nu_r=60.4$\,kHz and $T=17\,\mu$K). A hard region is the
special case $W=1/A_R$ inside and $0$ outside. On the uniform computation grid
the integral becomes the weighted sum
$\sum_i w_i\theta_i$ with $w_i=W_i/\sum_iW_i$, because the cell area cancels
between numerator and denominator; that is how the code evaluates it, on a
dedicated fine grid of $\pm4\sigma$, since the global grid is far too coarse for
$\sigma\approx0.1\,\mu$m.

\paragraph{What the atom does with the area.}
For a resonant two-level system with a real $\Omega(t)$ the Hamiltonian
commutes with itself at different times, so the excitation depends on the
accumulated area alone -- exactly, not approximately:
\begin{equation}
p_\uparrow(\mathbf r)=\sin^2\frac{\theta(\mathbf r)}{2},
\qquad
p_\uparrow(\mathbf r)=\frac{1}{1+\eta^2}
\sin^2\!\left(\sqrt{1+\eta^2}\,\frac{\theta(\mathbf r)}{2}\right)
\ \text{with the light shift},
\end{equation}
and the measured signal is $\int W p_\uparrow\,\mathrm d^2r$: $\sin^2$ is
non-linear, so the average is taken \emph{after} it. The spread of the area
itself, $U(\theta)=\sigma_W(\theta)/\langle\theta\rangle_W$, is what limits the
fidelity. Scattering multiplies the signal by $e^{-\Gamma_\mathrm{sc}T_p}$.

\paragraph{Working point.}
$3\times4$, Airy, waist $1.04\,\mu$m, $r_x/r_y=0.97/1.16$, all tone phases
zero, $\Delta=+50$\,GHz, $P=10\,\mu$W in the profile, $T_p=1\,\mu$s, atom
weighted ($\sigma=108$\,nm):

\begin{center}
\begin{tabular}{lr}
\hline
$\tilde P_\mathrm{prof}$ (simulation units) & $2.615\times10^{-11}$\,m$^2$\\
$\langle\tilde I\rangle_W$ (simulation units) & $1.913$\\
effective area $A_\mathrm{eff}$ & $13.67\,\mu$m$^2$\\
reference intensity $I_\mathrm{ref}=P/A_\mathrm{eff}$ & $73.2$\,W/cm$^2$\\
nominal Rabi frequency $f_\mathrm{Rabi}$ & $1.358$\,MHz\\
area at the reference, $2f_\mathrm{Rabi}T_p$ & $2.72\,\pi$\\
area at the best trigger, $\langle\theta\rangle_W$ & $11.54\,\pi$\\
mean Rabi frequency during that pulse & $5.77$\,MHz\\
$\sigma_\theta/\langle\theta\rangle_W$ there & $0.72\,\%$\\
power for a $\pi$ pulse of $1\,\mu$s at that trigger & $0.87\,\mu$W\\
\hline
\end{tabular}
\end{center}

Both routes -- the boxed formula and the detour over $f_\mathrm{Rabi}$ -- give
the same $11.539\,\pi$. The factor $4.25$ between the triggered and the
reference area is the rephasing peak of the twelve tones at phases zero.

\paragraph{Not included.}
Static two-photon detuning and magnetic field shifts, any uncompensated
gradient of the light shift, spontaneous emission during the pulse, the Zeeman
substructure, atomic motion during the pulse, finite pulse edges and the
acoustic fill time of the AOD, polarisation and vector light shifts,
intermodulation in the AOD and its diffraction efficiency. $P$ is the power
\emph{in the profile}, i.e. after both AODs.
